% part: intuitionistic-logic
% chapter: introduction
% section: bhk-interpretation

\documentclass[../../../include/open-logic-section]{subfiles}

\begin{document}

\olfileid{int}{int}{bhk}

\olsection{బ్రౌవర్--హైటింగ్--కొల్మొగొరోవ్ అర్థనిర్దేశం}

\begin{editorial}
  BHK అర్థనిర్దేశంతో అంతఃప్రజ్ఞావాద ప్రతిపాదనల
  చెల్లుబాటు నిరూపణలు అయోమయంగా ఉన్నాయి; వాటిని
  మెరుగ్గా వివరించాలి.
\end{editorial}

అంతఃప్రజ్ఞావాద సంధాయకాలకు ఒక అనౌపచారిక
నిర్మాణాత్మక అర్థనిర్దేశం ఉంది. దాన్ని సాధారణంగా
బ్రౌవర్--హైటింగ్--కొల్మొగొరోవ్ అర్థనిర్దేశం
అంటారు. ఇది ``నిర్మాణం'' అనే భావనను వాడుతుంది;
దాన్ని నిర్మాణాత్మక నిరూపణగా ఊహించవచ్చు.
(BHK అర్థనిర్దేశంలో ``నిరూపణ'' అనే పదాన్ని
వాడకపోవడానికి కారణం, ఔపచారిక
\tetoken{వ్యుత్పత్తి}{!!{derivation}} వ్యవస్థలోని
\tetoken{ఒక వ్యుత్పత్తి}{!!a{derivation}} భావనతో
గందరగోళం రాకూడదనే.) ఈ అంతఃప్రజ్ఞాత్మక భావన
ఆధారంగా BHK అర్థనిర్దేశం అంతఃప్రజ్ఞావాద
సంధాయకాల అర్థాలను వివరిస్తుంది.

\begin{enumerate}
\item అణు వాక్యానికి నిర్మాణం అంటే ఏమిటో మనకు
  తెలుసని అనుకుందాం.
\item $!A_1 \land !A_2$ నిర్మాణం ఒక జత
  $\tuple{M_1, M_2}$; ఇందులో $M_1$ అనేది~$!A_1$
  నిర్మాణం, $M_2$ అనేది~$!A_2$ నిర్మాణం.
\item $!A_1 \lor !A_2$ నిర్మాణం ఒక జత $\tuple{s, M}$;
  ఇందులో $s$ అనేది~$1$, $M$ అనేది~$!A_1$
  నిర్మాణం; లేదా $s$ అనేది~$2$, $M$ అనేది~$!A_2$
  నిర్మాణం.
\item $!A \lif !B$ నిర్మాణం,~$!A$ నిర్మాణాన్ని
  $!B$ నిర్మాణంగా మార్చే ప్రమేయకం.
\item $\lfalse$ (అసంభవం)కు నిర్మాణం లేదు.
\item $\lnot !A$ను $!A \lif \lfalse$కు పర్యాయంగా
  నిర్వచిస్తాం. అంటే $\lnot !A$ నిర్మాణం,
  $!A$ నిర్మాణాన్ని~$\lfalse$ నిర్మాణంగా మార్చే
  ప్రమేయకం.
\end{enumerate}

\begin{ex}
ఉదాహరణకు $\lnot \lfalse$ను తీసుకోండి. దాని
నిర్మాణం, $\lfalse$ యొక్క ఏ నిర్మాణాన్నైనా
ఇన్‌పుట్‌గా తీసుకుని~$\lfalse$ నిర్మాణాన్ని
అవుట్‌పుట్‌గా ఇచ్చే ప్రమేయకం. తాదాత్మ్య
ప్రమేయకం~$\Id{}$ అలాంటి నిర్మాణమే: $\lfalse$
నిర్మాణం~$M$ ఇస్తే, $\Id{}(M) = M$ ద్వారా
మళ్లీ~$\lfalse$ నిర్మాణమే లభిస్తుంది.
\end{ex}

సాధారణంగా $\lnot !A$ అంటే ``$!A$ నిర్మాణం
అసాధ్యం'' అని అర్థం.

\begin{ex}
ఏ ప్రతిపాదన~$!A$కైనా $!A \lif \lnot \lnot !A$ను
నిరూపిద్దాం; ఇది $!A \lif ((!A \lif \lfalse) \lif
\lfalse)$. నిర్మాణం ఒక ప్రమేయకం~$f$ అయి ఉండాలి:
$!A$ నిర్మాణం~$M$ ఇస్తే, అది $(!A \lif \lfalse)
\lif \lfalse$ నిర్మాణం~$f(M)$ను తిరిగి ఇస్తుంది.
$(!A \lif \lfalse) \lif
\lfalse$ నిర్మాణాన్ని~$f$ ఇలా నిర్మిస్తుంది:
$!A \lif \lfalse$ నిర్మాణం~$h$ను ఇన్‌పుట్‌గా
తీసుకుని~$\lfalse$ నిర్మాణాన్ని ఇచ్చే ప్రమేయకం
$g$ను నిర్వచించాలి. $g$ను ఇలా నిర్వచించవచ్చు:
ఇన్‌పుట్~$h$ను, మనకు ముందుగా లభించిన $!A$
నిర్మాణం~$M$కు వర్తింపజేయాలి. $h$ ఇచ్చే
ఫలితం~$h(M)$ ఒక $\lfalse$ నిర్మాణం కాబట్టి,
$M$ అనేది~$!A$ నిర్మాణమైతే $f(M)(h) = h(M)$
కూడా~$\lfalse$ నిర్మాణమే.
\end{ex}

\begin{ex}
$\lnot(!A \land \lnot !A)$కు, అంటే $(!A
\land (!A \lif \lfalse)) \lif \lfalse$కు,
ఒక నిర్మాణం ఇద్దాం. అది $!A \land (!A \lif
\lfalse)$ నిర్మాణం~$M$ను ఇన్‌పుట్‌గా తీసుకుని
$\lfalse$ నిర్మాణాన్ని ఇచ్చే ప్రమేయకం~$f$.
సంయోగం $!B_1 \land !B_2$ నిర్మాణం ఒక జత
$\tuple{N_1, N_2}$; ఇందులో $N_1$ అనేది~$!B_1$
నిర్మాణం, $N_2$ అనేది~$!B_2$ నిర్మాణం. $!B_1
\land !B_2$ నిర్మాణం నుంచి వరుసగా $!B_1$,
$!B_2$ నిర్మాణాలను తిరిగి పొందే ప్రమేయకాలు
$p_1$, $p_2$లను ఇలా నిర్వచించవచ్చు:
\begin{align*}
  p_1(\tuple{N_1, N_2}) &= N_1\\
  p_2(\tuple{N_1, N_2}) &= N_2
\end{align*}
$f$ ఇలా పనిచేస్తుంది: ముందుగా తన ఇన్‌పుట్~$M$కు
$p_1$ను వర్తింపజేసి~$!A$ నిర్మాణాన్ని పొందుతుంది.
తరువాత~$M$కు $p_2$ను వర్తింపజేసి $!A \lif
\lfalse$ నిర్మాణాన్ని పొందుతుంది. ఈ నిర్మాణం
ఒక ప్రమేయకం~$p_2(M)$; దానికి~$!A$ నిర్మాణం
ఇస్తే, అది~$\lfalse$ నిర్మాణాన్ని ఇస్తుంది.
అంటే $p_2(M)$ను $p_1(M)$కు వర్తింపజేస్తే
$\lfalse$ నిర్మాణం వస్తుంది. అందువల్ల
$f(M) = p_2(M)(p_1(M))$గా నిర్వచించవచ్చు.
\end{ex}

\begin{ex}
$((!A \land !B) \lif !C) \lif (!A \lif
(!B \lif !C))$కు నిర్మాణం ఇద్దాం. అంటే,
$(!A \land !B) \lif !C$ నిర్మాణం~$g$ను
$(!A \lif (!B \lif
!C))$ నిర్మాణంగా మార్చే ప్రమేయకం~$f$ కావాలి.
$g$ కూడా ఒక ప్రమేయకమే: $!A \land !B$
నిర్మాణాలను~$!C$ నిర్మాణాలుగా మారుస్తుంది.
\sourcecorrection{OLTEINTBHK-001}{మూలం ఈ ఒక్క చోట నిర్మాణపు ఫలితాన్ని $C$గా వ్రాసింది; ఇదే వాక్యంలోని మిగతా సూత్ర సంకేతానికి అనుగుణంగా $!C$గా సరిచేశాం.}
$f(g)$ అనే ఫలితం ఒక ప్రమేయకం~$h_g$;
అది $!A$ నిర్మాణాలను,~$!B$ నిర్మాణాల నుంచి
$!C$ నిర్మాణాలకు వెళ్లే ప్రమేయకాలుగా మారుస్తుంది.

సరే, ఇది అయోమయంగా ఉంది. ఇన్‌పుట్~$g$కు
$f$ ఇచ్చే అవుట్‌పుట్‌గా ఉండే ప్రమేయకం~$h_g$ను
నిర్మించాలి. $h_g$కు ఇన్‌పుట్,~$!A$
నిర్మాణం~$M$. $h_g(M)$ అవుట్‌పుట్,~$!B$
నిర్మాణం~$N$ను తీసుకుని~$!C$ నిర్మాణం ఇచ్చే
ప్రమేయకం~$k_M$ కావాలి. $k_{g,M}(N) =
g(\tuple{M, N})$గా పెట్టుకుందాం. $\tuple{M,
  N}$ అనేది~$!A \land !B$ నిర్మాణమని గుర్తుంచుకోండి.
కాబట్టి $k_{g,M}$ అనేది $!B \lif !C$ నిర్మాణం:
అది~$!B$ నిర్మాణం~$N$ను~$!C$ నిర్మాణంగా
మారుస్తుంది. ఇప్పుడు $h_g(M) = k_{g,M}$గా
పెట్టుకుందాం. ఈ ప్రమేయకం~$!A$ నిర్మాణం~$M$ను,
$!B \lif !C$ నిర్మాణం~$k_{g,M}$గా మారుస్తుంది.
ఇప్పుడు $f(g) = h_g$గా పెట్టుకుందాం. ఈ ప్రమేయకం
$(!A \land !B) \lif !C$ నిర్మాణం~$g$ను
$!A \lif (!B \lif !C)$ నిర్మాణంగా మారుస్తుంది.
అమ్మో!{}
\end{ex}

$!A \lor \lnot !A$ వాక్యాన్ని బహిష్కృత మధ్యమ
సూత్రం అంటారు. కొన్ని నిర్దిష్ట $!A$లకు
(ఉదాహరణకు $\lfalse \lor
\lnot \lfalse$కు) దీన్ని నిరూపించవచ్చు; కానీ
సాధారణంగా కాదు. ఎందుకంటే అంతఃప్రజ్ఞావాద
వియోజనంలో ఏదో ఒక వియోజ్యానికి నిర్మాణం కావాలి;
ప్రస్తుతం నిరూపించలేని, ఖండించలేని వాక్యాలు
(గోల్డ్‌బాక్ పరికల్పన వంటివి) ఉన్నాయి. అయితే
బహిష్కృత మధ్యమ సూత్రాన్ని ఖండించలేం కూడా:
అంటే $\lnot \lnot (!A
\lor \lnot !A)$ నిలుస్తుంది.

\begin{ex}
$\lnot \lnot (!A \lor \lnot !A)$ను నిరూపించడానికి,
$\lnot(!A \lor \lnot !A)$ నిర్మాణాన్ని, అంటే $(!A
\lor (!A \lif \lfalse)) \lif \lfalse$ నిర్మాణాన్ని,
$\lfalse$ నిర్మాణంగా మార్చే ప్రమేయకం~$f$ కావాలి.
అంటే $g$ అనేది $\lnot(!A
\lor \lnot !A)$ నిర్మాణమైతే $f(g)$ అనేది
$\lfalse$ నిర్మాణం అయ్యేలా $f$ను నిర్మించాలి.

$g$ అనేది $\lnot(!A \lor \lnot !A)$ నిర్మాణం,
అంటే $!A \lor \lnot !A$ నిర్మాణాన్ని
$\lfalse$ నిర్మాణంగా మార్చే ప్రమేయకం అనుకుందాం.
$!A \lor \lnot !A$ నిర్మాణం ఒక జత
$\tuple{s, M}$; ఇందులో $s = 1$ అయితే $M$
అనేది $!A$ నిర్మాణం, లేదా $s = 2$ అయితే
$M$ అనేది $\lnot !A$ నిర్మాణం.
$!A$ నిర్మాణం~$M_1$ను $!A \lor \lnot !A$
నిర్మాణంగా మార్చే ప్రమేయకం~$h_1$ అనుకుందాం:
అది $M_1$ను $\tuple{1, M_1}$గా మారుస్తుంది.
\sourcecorrection{OLTEINTBHK-002}{మూలం $h_1$కు వచ్చిన $M_1$ను రెండో భాగంలో $M_2$గా వ్రాసింది; అదే ఇన్‌పుట్‌తో జత కట్టేలా $M_1$గా సరిచేశాం.}
అలాగే $\lnot !A$ నిర్మాణం~$M_2$ను
$!A \lor \lnot !A$ నిర్మాణంగా మార్చే
ప్రమేయకం~$h_2$ అనుకుందాం: అది $M_2$ను
$\tuple{2, M_2}$గా మారుస్తుంది.

$k$ అనేది $\comp{h_1}{g}$ అనుకుందాం: $!A$
నిర్మాణం ఇస్తే~$\lfalse$ నిర్మాణాన్ని తిరిగి
ఇచ్చే ప్రమేయకం అది; అంటే $!A \lif \lfalse$,
లేదా $\lnot !A$, యొక్క నిర్మాణం. ఇప్పుడు
$l$ అనేది $\comp{h_2}{g}$ అనుకుందాం. దీనికి
$\lnot !A$ నిర్మాణం ఇస్తే~$\lfalse$
నిర్మాణాన్ని ఇస్తుంది. $k$ అనేది $\lnot !A$
నిర్మాణం కాబట్టి $l(k)$ అనేది~$\lfalse$
నిర్మాణం.

మొత్తంగా, $\lnot(!A \lor \lnot !A)$ నిర్మాణం~$g$ను
$\lfalse$ నిర్మాణంగా ఎలా మార్చవచ్చో వివరించాం.
అంటే $\lnot(!A \lor \lnot !A)$ నిర్మాణం~$g$ను
$\lfalse$ నిర్మాణం~$l(k)$గా మార్చే ప్రమేయకం~$f$,
$\lnot\lnot(!A \lor \lnot !A)$ నిర్మాణమే.
\end{ex}

మీరు చూసినట్లే, BHK అర్థనిర్దేశంతో
\tetoken{సూత్రాల}{!!{formula}s} అంతఃప్రజ్ఞావాద
చెల్లుబాటును చూపడం త్వరగా భారంగా,
అయోమయంగా మారుతుంది. అదృష్టవశాత్తూ,
అంతఃప్రజ్ఞావాద తర్కానికి మెరుగైన
\tetoken{వ్యుత్పత్తి}{!!{derivation}} వ్యవస్థలు,
మరింత ఖచ్చితమైన అర్థవిజ్ఞాన వ్యాఖ్యానాలు ఉన్నాయి.
\end{document}
